\section{Accounting and action-distribution properties}
\label{app:properties}

\subsection{Chronological passenger-time accounting}
\label{app:properties:ledger}

The following identity justifies the reward accounting in
\Cref{sec:problem:reward} and its relationship to the reported endpoint.

\begin{proposition}[Nonduplicating AEC ledger]
\label{prop:ledger}
Let $t_0$ be the first controllable batch time, and suppose the positive-duration
batch transitions partition $[t_0,T_{\mathrm{hor}}]$. Assume the passenger-count
functions are integrable. With the serialization in \Cref{sec:problem:smdp}
and costs in \Cref{eq:v7_cost},
\begin{equation}
\sum_{k=0}^{K-1}C_k
=\omega_{\mathrm{w}}\int_{t_0}^{T_{\mathrm{hor}}}N_{\mathrm{w}}(u)\,\mathrm{d}u
+\omega_{\mathrm{v}}\int_{t_0}^{T_{\mathrm{hor}}}N_{\mathrm{v}}(u)\,\mathrm{d}u.
\label{eq:ledger_identity}
\end{equation}
Define the pre-control cost
\begin{equation}
C_{\mathrm{pre}}=
\omega_{\mathrm{w}}\int_0^{t_0}N_{\mathrm{w}}(u)\,\mathrm{d}u
+\omega_{\mathrm{v}}\int_0^{t_0}N_{\mathrm{v}}(u)\,\mathrm{d}u.
\label{eq:precontrol_cost}
\end{equation}
For $N_{\mathrm{dep}}>0$, the primary endpoint satisfies
\begin{equation}
Y=\frac{C_{\mathrm{pre}}+\sum_{k=0}^{K-1}C_k}{N_{\mathrm{dep}}}.
\label{eq:endpoint_ledger}
\end{equation}
For a fixed exogenous realization, $C_{\mathrm{pre}}$ and $N_{\mathrm{dep}}$
are policy-invariant under the registered demand contract.
\end{proposition}

\begin{proof}
Every nonfinal token in a simultaneous batch has $\Delta t_k=0$, so both
integrals defining its cost vanish. The final token carries exactly the
interval from that batch to the next physical batch or the horizon.
These intervals do not overlap except at endpoints of zero measure. Finite
additivity of the integrals therefore gives \Cref{eq:ledger_identity}.
Adding the initial interval in \Cref{eq:precontrol_cost} gives the full-horizon
integrals in \Cref{eq:completion_rate}, proving \Cref{eq:endpoint_ledger}.
There is no control before $t_0$, and passenger departures are fixed by the
exogenous demand realization, which proves the final statement.
\end{proof}

\newpage
\subsection{Hurdle density and deterministic action}
\label{app:properties:hurdle}

\Cref{sec:method:architecture} uses the following distributional property to
compute likelihoods and specify deterministic evaluation.

\begin{proposition}[Normalized hurdle policy and positive-branch median]
\label{prop:hurdle}
For finite $g_{\phi}(s)$ and $\mu_{\phi}(s)$ and
$\sigma_{\phi}(s)>0$, let $\mathfrak{n}_{\phi}(\cdot\mid s)$ denote the
normal density with mean $\mu_{\phi}(s)$ and standard deviation
$\sigma_{\phi}(s)$. The conditional positive-branch density is
\begin{equation}
f_{\phi}(a\mid s)=
\frac{2\mathfrak{n}_{\phi}(\zeta(a)\mid s)}{1-a^2},
\qquad
\zeta(a)=\log\frac{1+a}{1-a},\quad -1<a<1,
\label{eq:positive_density}
\end{equation}
where $\zeta(a)$ is the latent coordinate associated with action $a$.
The policy in \Cref{eq:hurdle_policy} is a probability measure, and its
conditional positive-branch median is $2\sigmoid(\mu_{\phi}(s))-1$.
\end{proposition}

\begin{proof}
The map $a=2\sigmoid(\zeta)-1$ is a strictly increasing bijection from
$\R$ to $(-1,1)$, with derivative
$\mathrm{d}a/\mathrm{d}\zeta=(1-a^2)/2$. Change of variables gives
\Cref{eq:positive_density} and
$\int_{-1}^{1}f_{\phi}(a\mid s)\,\mathrm{d}a=1$.
The continuous mass is therefore $p^+_{\phi}(s)$; adding the point mass
$1-p^+_{\phi}(s)$ yields total mass one. The normal latent variable has median
$\mu_{\phi}(s)$, and a strictly increasing transformation preserves its
median. This proves the deterministic-action statement; nonlinear
transformation does not generally preserve the conditional mean.
\end{proof}

